\chapter{AI-Ready Data}
\label{ch:aireadydata}

For decades, people working in biodiversity have published occurrence records,
sequences, taxonomies, traits and time series openly. They curated, standardised
and documented these data so that others could use them and add to the
knowledge of ecosystems.

\section{What the term hides}

Data is called \term{AI-ready} when it is machine-readable, uses controlled
vocabularies, carries structured metadata, can be downloaded in bulk and is
openly licensed. Every one of those is good data management, which
biodiversity informatics was doing long before the term existed. What the term
adds is a consumer --- often a company that will ingest the corpus, build a
model, and return neither the weights nor credit to the curators. This is not an
argument against openness, which has served ecology well. It is an argument for
reading the terms.

\begin{keypoint}
When someone asks you to make data ``AI-ready'', ask: ready for whom, under what
licence, with what attribution, and does anything come back to the people who
maintain it?
\end{keypoint}

\section{Standards}

\textbf{Describing data.} Darwin Core is the backbone of occurrence exchange,
packaged as Darwin Core Archives. ABCD covers collections, while MIxS sets the
minimum information for sequence data. These descriptions let us compare a
metagenome from a Greek coast with one from the North Sea. Ontologies do the semantic work: ENVO
for environments, EUNIS for habitats, the Gene Ontology and others. Choosing
between overlapping vocabularies --- ENVO, EUNIS or LUCAS for a habitat field ---
is a real design decision. I-ADOPT breaks a measured variable into its parts so
that ``sea surface temperature'' from two networks can be recognised as the same
thing \parencite{magagna2021iadopt}, and DataHarmonizer validates messy
contextual metadata against a specification
\parencite{gill2023dataharmonizer}. This work makes it possible to combine data
collected by different people for different purposes.

The microbiome comparison by \textcite{venetsianou2026relatedness} shows why
descriptions need as much care as the data files. Generic descriptions and
abstracts reused across many studies could obscure differences between them;
the authors filtered heavily reused abstracts before comparing study pairs.
The model could only assess the context supplied to it. For a reusable
microbiome dataset, I would therefore retain the sampling conditions, links
between samples and studies, and the history of laboratory and computational
processing alongside the sequences.

Jonathan Jacobs and Patrick Boyle make a related argument in their Undark
opinion essay on AI-driven biology, listed in \Cref{app:tutorials}. They stress
traceable biological materials and records of handling and origin. Their
argument extends provenance beyond a downloadable table to the material that
produced it. My reading for biodiversity work is that making records easy to
retrieve cannot compensate for losing their connection to specimens, cultures
or field observations.

\textbf{Machine interfaces.} POSIX, the Unix file and process conventions,
underlies everything \parencite{ritchie1974unix}. OpenAPI describes web APIs so
that an agent can discover and call them. The \term{Model Context Protocol}
(MCP) is a standard way to offer tools and data to language models: a resource is
wrapped once instead of once per client. EMBL-EBI's InterPro server is a good
example. An agent can ask in plain language which domains a protein has or
which structures contain a family. The server calls the right InterPro services,
with no account needed \parencite{interpro2026mcp}. SPARQL and RDF serve
knowledge graphs, which models can now query \parencite{smeros2025sparqlllm}.
Frameworks such as LangChain and LangGraph chain these pieces into applications,
and Ollama runs open models locally (\Cref{ch:stack}).

\section{Infrastructure}

GBIF aggregates occurrence records worldwide and OBIS does the same for the sea.
PlutoF manages data from collection to publication, with UNITE for fungal
sequences \parencite{koljalg2019plutof}; BOLD holds barcodes and ENA, GenBank
and SRA the rest. These are long-lived and publicly funded. Their survival is a
policy achievement, not a market outcome, and it is not guaranteed. TDWG governs
Darwin Core, ELIXIR coordinates European life-science data and names the Core
Data Resources the community commits to sustaining, and the LF AI \& Data
Foundation hosts much open model tooling. Open standards outlive the
organisations that introduce them \parencite{wilson2013foss}.

Hugging Face has become the default place to share open models and datasets. It
is convenient, a single commercial point of failure, and an attack surface:
downloading a model can mean running someone else's code
\parencite{openai2026huggingface}. Prefer the \code{safetensors} format and
check checksums.

\section{Open weights and open source}

\textbf{Open weight} means you can download the model's parameters, run it and
fine-tune it --- Gemma, Kimi, Qwen and GPT-OSS models are examples. It does not
mean you can see what it was trained on or rebuild it. \textbf{Open source},
properly, would include the training data, code and pipeline as well. Almost
every model sold as open is only open-weight. \textcite{hai2026openweight} argue
this is not enough for science: without the training data you cannot check
whether your test set was in it, audit bias, or reproduce the result. Licences
vary too; Apache-2.0 and MIT are genuinely open, while several ``open'' model
licences restrict use.

\begin{table}[htbp]
\centering
\caption{Degrees of openness. Only the last row supports reproduction.}
\label{tab:openness}
\small
\begin{tabularx}{\textwidth}{lXX}
\toprule
\textbf{Level} & \textbf{You get} & \textbf{You cannot} \\
\midrule
API only      & An endpoint & Run offline; pin a version; audit anything \\
Open weight, restricted licence & Weights, conditional use & Use freely; audit training data \\
Open weight, Apache/MIT & Weights, unrestricted use & Audit training data; reproduce \\
Fully open source & Weights, data, code, pipeline & --- \\
\bottomrule
\end{tabularx}
\end{table}

\section{Retrieval-augmented generation}

RAG gives a model knowledge it was not trained on, without retraining. Ask a
model ``what colour is my ball?'' and it will invent an answer. Say ``I have a
green box with a red ball in it --- what colour is my ball?'' and it answers
correctly. RAG automates that second clause: a search system finds relevant
passages and adds them to the prompt. The steps:

\begin{enumerate}
  \item \textbf{Convert} documents to text. For PDFs this is the hardest step
    and limits everything after it; Docling, which understands page layout and
    tables, is much better than plain extraction \parencite{auer2024docling}.
  \item \textbf{Chunk} into passages, respecting the document's structure.
  \item \textbf{Embed} each chunk as a vector (\code{bge-m3}, \code{nomic-embed-text}).
  \item \textbf{Index} them, for example in PostgreSQL with \code{pgvector}.
  \item \textbf{Retrieve} with keyword and vector search combined, then
    \term{rerank} with a more careful model such as \code{bge-reranker}.
  \item \textbf{Generate} an answer that must cite the retrieved passages.
\end{enumerate}

\textcite{kurniawan2024ragfailures} list seven ways it fails. Three are retrieval
failures: the answer is not in the corpus, it is ranked too low, or it is cut
from the context. Four are generation failures: the model misses it, uses the
wrong format, is too general or too narrow, or gives an incomplete answer. Log
what was retrieved separately from what was generated, or you cannot tell which
half broke. For a small, stable corpus, a curated set of notes placed directly
in the context can beat retrieval altogether; test before assuming you need RAG.
